% =============================================================================
% Module 1b: Differential Geometry and the Metric Tensor
% =============================================================================
% Sources: Carroll Ch. 2--3 (Secs. 2.1--2.4, 3.1--3.4),
%          Congdon & Keeton Ch. 3.2
% Mathematica: ../Mathematica/01b_Differential_Geometry/
% =============================================================================

\chapter{Differential Geometry and the Metric Tensor}
\label{ch:differential_geometry}

\begin{quote}
\textit{The appropriate mathematical structure used to describe curvature is
that of a differentiable manifold: essentially, a kind of set that looks
locally like flat space, but might have a very different global geometry.}
\hfill --- Sean Carroll, \textit{Spacetime and Geometry}
\end{quote}

\noindent
In Module~1a we worked exclusively in flat Minkowski spacetime.  To
describe gravity, we must allow spacetime to be \emph{curved}.  This
chapter introduces the mathematical machinery of differential geometry
--- manifolds, the general metric tensor, the covariant derivative,
Christoffel symbols, and the curvature tensors --- that forms the
foundation of general relativity and, ultimately, gravitational lensing.

\medskip
\noindent
\textbf{Companion Mathematica files:}
\begin{itemize}[nosep]
    \item \mathlink{01b_Differential_Geometry/christoffel_symbols.wl}{christoffel\_symbols.wl}
          --- general-purpose symbolic computation of Christoffel symbols
          and curvature tensors from any input metric
    \item \mathlink{01b_Differential_Geometry/curved_surfaces.nb}{curved\_surfaces.nb}
          --- interactive visualizations of parallel transport, geodesics,
          and curvature on 2D surfaces
\end{itemize}

% =====================================================================
\section{Manifolds and Coordinate Charts}
\label{sec:manifolds}

A \textbf{manifold} is a space that may be curved globally but looks
like $\mathbb{R}^n$ in every sufficiently small neighborhood.  More
precisely, an $n$-dimensional manifold $M$ is a topological space
equipped with a collection of \textbf{coordinate charts} --- invertible
maps $\phi : U \to \mathbb{R}^n$ from open subsets $U \subset M$ ---
that cover all of $M$ and are smoothly compatible on overlaps
(Carroll Sec.~2.2).

Examples of manifolds include:
\begin{itemize}[nosep]
    \item $\mathbb{R}^n$ (trivially)
    \item The 2-sphere $S^2$ (the surface of a ball)
    \item The 2-torus $T^2$ (the surface of a doughnut)
    \item Spacetime itself (a 4-dimensional Lorentzian manifold)
\end{itemize}

A key fact: in general, \textit{no single coordinate system can cover an
entire manifold}.  The sphere $S^2$, for instance, requires at least two
charts (e.g., stereographic projections from the north and south poles;
Carroll eqs.~2.9--2.10).  We will always work locally in a single chart,
keeping in mind that our results are coordinate-independent.


% =====================================================================
\section{The General Metric Tensor}
\label{sec:general_metric}

On a manifold, we need a way to measure distances and angles.  This is
provided by the \textbf{metric tensor} $g_{\mu\nu}(x)$, a symmetric,
nondegenerate $(0,2)$ tensor field.  The line element is
(cf.\ Carroll eq.~1.20, now generalized):
\begin{equation}
    \boxed{
    ds^2 = g_{\mu\nu}(x)\, dx^\mu\, dx^\nu
    }
    \label{eq:general_line_element}
\end{equation}
In flat Minkowski space, $g_{\mu\nu} = \eta_{\mu\nu}$ everywhere.  In
curved spacetime, $g_{\mu\nu}$ varies from point to point and encodes
the gravitational field.

The \textbf{inverse metric} $g^{\mu\nu}$ is defined by
\begin{equation}
    g^{\mu\lambda} g_{\lambda\nu} = \delta^\mu_{\ \nu} \,.
    \label{eq:inverse_metric}
\end{equation}
As before, the metric raises and lowers indices:
$V_\mu = g_{\mu\nu} V^\nu$, $V^\mu = g^{\mu\nu} V_\nu$.

\begin{example}
The metric on the 2-sphere of radius $R$ in spherical coordinates
$(\theta, \phi)$ is
\begin{equation}
    ds^2 = R^2\bigl(d\theta^2 + \sin^2\theta\, d\phi^2\bigr) \,,
    \label{eq:sphere_metric}
\end{equation}
so $g_{\theta\theta} = R^2$, $g_{\phi\phi} = R^2\sin^2\theta$, and all
other components vanish.
\end{example}

\begin{example}
\label{ex:polar_metric}
The flat plane in polar coordinates $(r, \phi)$ has metric
\begin{equation}
    ds^2 = dr^2 + r^2\, d\phi^2 \,.
    \label{eq:polar_metric}
\end{equation}
Despite the coordinate-dependent appearance, this space is flat ---
the curvature tensors vanish.  We will verify this in Mathematica.
\end{example}


% =====================================================================
\section{Vectors and Tensors on Manifolds}
\label{sec:vectors_manifolds}

On a general manifold, vectors live in the \textbf{tangent space} $T_p$
at each point $p$.  A vector $V$ at $p$ is identified with a directional
derivative operator along a curve $x^\mu(\lambda)$ through $p$
(Carroll Sec.~2.3):
\begin{equation}
    V = V^\mu \partial_\mu \,,
    \label{eq:vector_as_derivative}
\end{equation}
where $\partial_\mu \equiv \partial/\partial x^\mu$ form a
\textbf{coordinate basis} for $T_p$.

Under a coordinate transformation $x^\mu \to x^{\mu'}$, vector
components transform as (Carroll eq.~2.19):
\begin{equation}
    V^{\mu'} = \frac{\partial x^{\mu'}}{\partial x^\nu}\, V^\nu \,.
    \label{eq:vector_transform_general}
\end{equation}
This generalizes the Lorentz transformation law from Module~1a to
arbitrary coordinate changes.

Dual vectors (one-forms) and general tensors transform analogously, with
one factor of $\partial x^{\mu'}/\partial x^\nu$ for each upper index
and $\partial x^\nu/\partial x^{\mu'}$ for each lower index.


% =====================================================================
\section{The Covariant Derivative}
\label{sec:covariant_derivative}

Ordinary partial derivatives of tensors are \textit{not} tensors on a
curved manifold: the transformation of $\partial_\mu V^\nu$ picks up an
extra, non-tensorial term from the coordinate dependence of the basis
vectors (Carroll Sec.~3.2).

To fix this, we introduce the \textbf{covariant derivative} $\nabla_\mu$,
which adds a correction term involving the \textbf{connection coefficients}
$\Gamma^\rho_{\ \mu\sigma}$ (Carroll eq.~3.5):
\begin{equation}
    \boxed{
    \nabla_\mu V^\nu = \partial_\mu V^\nu + \Gamma^\nu_{\ \mu\lambda}\, V^\lambda
    }
    \label{eq:cov_deriv_vector}
\end{equation}
For a one-form (Carroll eq.~3.16):
\begin{equation}
    \nabla_\mu \omega_\nu = \partial_\mu \omega_\nu
    - \Gamma^\lambda_{\ \mu\nu}\, \omega_\lambda \,.
    \label{eq:cov_deriv_oneform}
\end{equation}
The general rule for a tensor of arbitrary rank: $+\Gamma$ for each upper
index, $-\Gamma$ for each lower index (Carroll eq.~3.17).

The connection coefficients are \textit{not} tensors --- they transform
inhomogeneously (Carroll eq.~3.10) --- but the covariant derivative
$\nabla_\mu V^\nu$ \textit{is} a tensor.


% =====================================================================
\section{Christoffel Symbols}
\label{sec:christoffel}

In GR we demand two additional properties of the connection:
\begin{enumerate}[nosep]
    \item \textbf{Torsion-free:}
          $\Gamma^\lambda_{\ \mu\nu} = \Gamma^\lambda_{\ \nu\mu}$
    \item \textbf{Metric compatibility:}
          $\nabla_\rho g_{\mu\nu} = 0$
\end{enumerate}
These uniquely determine the \textbf{Christoffel symbols} (also called
the Levi-Civita connection; Carroll eq.~3.27):
\begin{equation}
    \boxed{
    \Gamma^\sigma_{\ \mu\nu} = \frac{1}{2}\, g^{\sigma\rho}
    \bigl(\partial_\mu g_{\nu\rho} + \partial_\nu g_{\rho\mu}
    - \partial_\rho g_{\mu\nu}\bigr)
    }
    \label{eq:christoffel}
\end{equation}
This is one of the most important formulas in GR.  Given any metric,
we can compute the Christoffel symbols and from them all of the geometry.

\begin{example}
For the plane in polar coordinates, $ds^2 = dr^2 + r^2 d\phi^2$
(Example~\ref{ex:polar_metric}), the only nonvanishing Christoffel
symbols are (Carroll eqs.~3.29--3.31):
\begin{equation}
    \Gamma^r_{\ \phi\phi} = -r \,, \qquad
    \Gamma^\phi_{\ r\phi} = \Gamma^\phi_{\ \phi r} = \frac{1}{r} \,.
    \label{eq:christoffel_polar}
\end{equation}
We derive these and more general cases in
\mathlink{01b_Differential_Geometry/christoffel_symbols.wl}{christoffel\_symbols.wl}.
\end{example}


% =====================================================================
\section{Parallel Transport}
\label{sec:parallel_transport}

In flat space, we can freely slide a vector from one point to another
while keeping it ``constant.''  On a curved manifold, this is no longer
unambiguous --- the result depends on the \textit{path taken}.

The notion of keeping a vector constant along a curve $x^\mu(\lambda)$
is formalized as \textbf{parallel transport}: a vector $V^\mu$ is
parallel-transported if (Carroll eq.~3.40)
\begin{equation}
    \frac{d V^\mu}{d\lambda}
    + \Gamma^\mu_{\ \rho\sigma}\, \frac{dx^\rho}{d\lambda}\, V^\sigma
    = 0 \,.
    \label{eq:parallel_transport}
\end{equation}
An important consequence of metric compatibility is that parallel
transport preserves the inner product: if $V^\mu$ and $W^\mu$ are both
parallel-transported along a curve, then
$g_{\mu\nu} V^\mu W^\nu = \text{const}$ (Carroll eq.~3.42).

On the 2-sphere, parallel transporting a vector around a closed loop
results in a rotation of the vector --- the angle of rotation is
proportional to the enclosed area and the curvature.  This is visualized
interactively in
\mathlink{01b_Differential_Geometry/curved_surfaces.nb}{curved\_surfaces.nb}.


% =====================================================================
\section{The Geodesic Equation}
\label{sec:geodesic_equation}

A \textbf{geodesic} is the curved-space generalization of a straight line:
it is the path that parallel-transports its own tangent vector.
Setting $V^\mu = dx^\mu/d\lambda$ in~\eqref{eq:parallel_transport}, we
obtain the \textbf{geodesic equation} (Carroll eq.~3.44):
\begin{equation}
    \boxed{
    \frac{d^2 x^\mu}{d\lambda^2}
    + \Gamma^\mu_{\ \rho\sigma}\,
    \frac{dx^\rho}{d\lambda}\, \frac{dx^\sigma}{d\lambda} = 0
    }
    \label{eq:geodesic}
\end{equation}
This can equivalently be derived by extremizing the proper time/interval
functional (Carroll Sec.~3.3, eqs.~3.45--3.55; Congdon \& Keeton
eq.~3.60).

Key properties:
\begin{itemize}
    \item \textbf{Timelike geodesics} ($ds^2 < 0$): paths of freely falling
          massive particles; parameterized by proper time $\tau$.
    \item \textbf{Null geodesics} ($ds^2 = 0$): paths of light rays;
          parameterized by an affine parameter $\lambda$.
    \item \textbf{Spacelike geodesics} ($ds^2 > 0$): shortest spatial
          distances.
    \item The character (timelike/null/spacelike) is preserved along the
          geodesic (Carroll Sec.~3.4).
\end{itemize}

In flat space with Cartesian coordinates, $\Gamma^\mu_{\ \rho\sigma} = 0$,
so the geodesic equation reduces to $d^2 x^\mu/d\lambda^2 = 0$ --- a
straight line, as expected.


% =====================================================================
\section{Curvature Tensors}
\label{sec:curvature}

The \textbf{Riemann curvature tensor} $R^\rho_{\ \sigma\mu\nu}$ measures
how much parallel transport around a closed loop rotates a vector.  It is
defined by (Carroll eq.~3.4):
\begin{equation}
    \boxed{
    R^\rho_{\ \sigma\mu\nu} = \partial_\mu \Gamma^\rho_{\ \nu\sigma}
    - \partial_\nu \Gamma^\rho_{\ \mu\sigma}
    + \Gamma^\rho_{\ \mu\lambda}\, \Gamma^\lambda_{\ \nu\sigma}
    - \Gamma^\rho_{\ \nu\lambda}\, \Gamma^\lambda_{\ \mu\sigma}
    }
    \label{eq:riemann}
\end{equation}

\textbf{A manifold is flat if and only if $R^\rho_{\ \sigma\mu\nu} = 0$
everywhere.}  This is the mathematical statement that curvature is
characterized entirely by the Riemann tensor.

From the Riemann tensor we construct:
\begin{itemize}
    \item \textbf{Ricci tensor} (contraction on 1st and 3rd indices):
    \begin{equation}
        R_{\mu\nu} = R^\lambda_{\ \mu\lambda\nu}
        \label{eq:ricci_tensor}
    \end{equation}
    \item \textbf{Ricci scalar} (trace of the Ricci tensor):
    \begin{equation}
        R = g^{\mu\nu} R_{\mu\nu}
        \label{eq:ricci_scalar}
    \end{equation}
\end{itemize}

These are the quantities that appear in Einstein's field equation
\begin{equation}
    R_{\mu\nu} - \frac{1}{2} R\, g_{\mu\nu} = 8\pi G\, T_{\mu\nu} \,,
    \label{eq:einstein_field_eq}
\end{equation}
relating the curvature of spacetime (left side) to its matter and energy
content (right side).

\begin{remark}
The Riemann tensor in $n$ dimensions has $n^2(n^2 - 1)/12$ independent
components.  In 4D spacetime, that is 20.  In 2D, there is only 1 ---
the Gaussian curvature determines everything.
\end{remark}


% =====================================================================
\section{The Metric Determines Everything}
\label{sec:metric_determines}

Let us collect the chain of logic:
\begin{center}
\textbf{Metric} $g_{\mu\nu}$ \\
$\downarrow$ eq.~\eqref{eq:christoffel} \\
\textbf{Christoffel symbols} $\Gamma^\sigma_{\ \mu\nu}$ \\
$\downarrow$ eq.~\eqref{eq:geodesic} \qquad $\downarrow$ eq.~\eqref{eq:riemann} \\
\textbf{Geodesics} (particle paths) \qquad \textbf{Curvature} $R^\rho_{\ \sigma\mu\nu}$
\end{center}

\noindent
This is the central message: \textit{the metric is the fundamental object
in GR}.  Once you have $g_{\mu\nu}$, you can compute everything ---
how particles move, how light bends, and how spacetime curves.

For gravitational lensing, the key application is: given the metric
around a massive object (e.g., the Schwarzschild metric for a point mass),
solve the null geodesic equation to trace light rays and determine how
images are formed.


% =====================================================================
\section{Summary}
\label{sec:dg_summary}

\begin{itemize}
    \item A manifold is locally flat but may have global curvature; the
          metric $g_{\mu\nu}(x)$ encodes the geometry.
    \item The covariant derivative $\nabla_\mu$ generalizes the partial
          derivative to curved spacetime, using the connection
          $\Gamma^\sigma_{\ \mu\nu}$.
    \item The Christoffel symbols are uniquely determined by the metric
          via~\eqref{eq:christoffel} (torsion-free, metric-compatible).
    \item Parallel transport depends on the path; this path-dependence is
          measured by the Riemann curvature tensor.
    \item Geodesics --- paths that parallel-transport their own tangent
          vector --- are the trajectories of freely falling particles
          and light rays.
    \item Einstein's equation relates curvature ($R_{\mu\nu}$) to
          matter content ($T_{\mu\nu}$).
\end{itemize}


% =====================================================================
\section{Problems}
\label{sec:dg_problems}

\begin{exercise}
\label{ex:christoffel_sphere}
Compute all nonvanishing Christoffel symbols for the 2-sphere metric
$ds^2 = R^2(d\theta^2 + \sin^2\theta\, d\phi^2)$ by hand, then verify
with \mathlink{01b_Differential_Geometry/christoffel_symbols.wl}{christoffel\_symbols.wl}.
\end{exercise}

\begin{exercise}
\label{ex:christoffel_polar}
For the flat plane in polar coordinates, $ds^2 = dr^2 + r^2 d\phi^2$:
\begin{enumerate}[label=(\alph*)]
    \item Compute the Christoffel symbols.
    \item Compute the Riemann tensor and verify that the space is flat
          ($R^\rho_{\ \sigma\mu\nu} = 0$).
    \item Write down the geodesic equations and show that straight lines
          in Cartesian coordinates ($x = a + b\lambda$, $y = c + d\lambda$)
          satisfy them.
\end{enumerate}
\end{exercise}

\begin{exercise}
\label{ex:ricci_sphere}
For the 2-sphere metric, compute:
\begin{enumerate}[label=(\alph*)]
    \item The Riemann tensor $R^\theta_{\ \phi\theta\phi}$.
    \item The Ricci tensor $R_{\theta\theta}$ and $R_{\phi\phi}$.
    \item The Ricci scalar $R$.  Show that $R = 2/R^2$ (constant
          positive curvature, as expected for a sphere of radius $R$).
\end{enumerate}
\end{exercise}

\begin{exercise}
\label{ex:parallel_transport_sphere}
A vector initially pointing east is parallel-transported from the equator
of a unit sphere ($R = 1$) along the following path: north along the
meridian $\phi = 0$ to the pole, then south along $\phi = \alpha$ back to
the equator.  Show that the vector has rotated by angle $\alpha$.
(Hint: solve the parallel transport equation~\eqref{eq:parallel_transport}
along each segment, or use the geometric argument that the rotation equals
the solid angle enclosed.)
\end{exercise}

\begin{exercise}
\label{ex:geodesic_sphere}
Show that great circles are geodesics of the 2-sphere.
Specifically, show that the equator $\theta = \pi/2$, $\phi = \lambda$
satisfies the geodesic equation~\eqref{eq:geodesic} with the Christoffel
symbols from Exercise~\ref{ex:christoffel_sphere}.
\end{exercise}
